% Brief (tikz-diagram-craft). Register row ch7-12 (should).
% Reader question: which rules can be prevented before the fact, and which
%   are only caught after?
% Claim: Synchronous State Decidability is the exact boundary. A rule
%   decidable against deterministic local DB state and the transaction
%   payload is regimented (its violation cannot occur at all); a rule
%   needing async checks, external I/O, or non-deterministic oracles is
%   relegated to post-commit enforcement (violation occurs, then is caught
%   and salvaged).
% Evidence: website-v2/public/whitepaper/agent-transactions-whitepaper.tex
%   OP-2 paragraph (Regimentation vs. Enforcement), citing Jones and Sergot
%   1993's regimentation/enforcement distinction; the three named compiled
%   invariants NoteMonotonicity, EscrowInvariant, LockOwnerValid map
%   directly to sec:tlaplus's safety properties; Arbiter-detected
%   violations are the enforced side, triggering the salvage protocol.
% Must distinguish: which of the two mechanisms is structurally impossible
%   to violate (regimented) from which is possible-but-caught (enforced).
% Grammar chosen: two-cell boundary, mirroring
%   whitepaper/figures/fig-swk-controllability-quadrant.tex's grammar (a
%   shaded marked cell, one concrete rule set per cell) collapsed to the
%   one real axis this rule actually has -- Synchronous State Decidability
%   -- rather than forcing a second, unsupported axis into a fake 2x2.
%   Rejected: a literal four-cell quadrant with sync-decidability crossed
%   against pre/post-commit, which the chapter's own text states are the
%   SAME distinction restated, not two independent axes; two empty or
%   duplicate cells would violate S8's no-empty-box rule.
% Five-second test: a reader sees three named invariants on the left
%   (structurally impossible to violate) and the Arbiter on the right
%   (violation possible, then caught).
\begin{figure}[H]
\centering
\pdfigurehue{pdgold}%
\begin{tikzpicture}[pd figure,x=1cm,y=1cm,
  cell/.style={font=\pdfiglabel,align=center,text width=4.35cm}]
  \fill[pd focus fill] (0,0) rectangle (4.55,4.6);
  \draw[pd rule] (0,0) rectangle (9.1,4.6);
  \draw[pd spine] (4.55,0) -- (4.55,4.6);
  \node[pd title,anchor=south] at (2.28,4.75) {regimented (pre-commit)};
  \node[pd title,anchor=south] at (6.83,4.75) {enforced (post-commit)};
  \node[pd axis label,anchor=north,text width=8.6cm,align=center] at (4.55,-.25)
    {Synchronous State Decidability: decidable against local DB state $\Rightarrow$ regimented; needs async I/O or an oracle $\Rightarrow$ enforced};
  \node[cell] at (2.28,2.9) {$\bigstar$ structurally impossible to violate};
  \node[cell] at (2.28,1.5) {\texttt{NoteMonotonicity}, \texttt{EscrowInvariant}, \texttt{LockOwnerValid}};
  \node[cell] at (6.83,2.9) {violation occurs, then caught};
  \node[cell] at (6.83,1.5) {Arbiter-detected violations $\to$ salvage protocol};
\end{tikzpicture}
\caption{Regimentation versus enforcement, in the sense of Jones and
Sergot's normative-systems distinction: a rule falls on one side or the
other of Synchronous State Decidability, never both. The three compiled
invariants (left, shaded, marked $\bigstar$) are checked against
deterministic local state and the transaction payload alone, so their
violation cannot occur at all. Every other invariant needs an async check,
external I/O, or a non-deterministic oracle, so the Arbiter (right) can
only detect a violation after the fact and route it to the salvage
protocol. [internal]}
\label{fig:bonded-regiment-boundary}
\end{figure}
